\documentclass[11pt,a4paper]{article}
\usepackage[margin=2.5cm,headheight=15pt]{geometry}
\usepackage{amsmath,amssymb,amsthm}
\usepackage{mathtools}
\usepackage{microtype}
\usepackage{hyperref}
\usepackage{cleveref}
\usepackage{enumitem}
\usepackage{mdframed}
\usepackage{xcolor}
\usepackage{titlesec}
\usepackage{booktabs}
\usepackage{fancyhdr}

\definecolor{navyblue}{RGB}{26,58,92}
\definecolor{midblue}{RGB}{44,95,138}
\definecolor{lightblue}{RGB}{232,240,248}
\definecolor{proofgray}{RGB}{242,242,242}
\definecolor{opencolor}{RGB}{235,245,235}
\definecolor{openborder}{RGB}{0,120,0}
\definecolor{warncolor}{RGB}{255,245,220}
\definecolor{warnborder}{RGB}{180,100,0}

\hypersetup{colorlinks=true,linkcolor=midblue,citecolor=midblue,urlcolor=midblue,
  pdftitle={SSP Part X: Position-Indexed Formulations}}

\pagestyle{fancy}\fancyhf{}
\fancyhead[L]{\small\color{navyblue}\textit{Part X: Position-Indexed Formulations of the Uniform SSP}}
\fancyhead[R]{\small\color{navyblue}\thepage}
\renewcommand{\headrulewidth}{0.4pt}

\titleformat{\section}{\large\bfseries\color{navyblue}}{}{0em}{}[\color{navyblue}\titlerule]
\titleformat{\subsection}{\normalsize\bfseries\color{midblue}}{}{0em}{}

\newmdenv[backgroundcolor=lightblue,linecolor=navyblue,linewidth=1.2pt,
  leftline=true,rightline=false,topline=false,bottomline=false,
  innerleftmargin=12pt,innerrightmargin=8pt,innertopmargin=8pt,
  innerbottommargin=8pt,skipabove=10pt,skipbelow=6pt]{thmbox}
\newmdenv[backgroundcolor=proofgray,linecolor=midblue,linewidth=0.8pt,
  leftline=true,rightline=false,topline=false,bottomline=false,
  innerleftmargin=10pt,innerrightmargin=8pt,innertopmargin=6pt,
  innerbottommargin=6pt,skipabove=6pt,skipbelow=6pt]{proofbox}
\newmdenv[backgroundcolor=lightblue,linecolor=navyblue,linewidth=1.5pt,roundcorner=3pt,
  innerleftmargin=14pt,innerrightmargin=14pt,innertopmargin=10pt,innerbottommargin=10pt,
  skipabove=12pt,skipbelow=8pt]{formbox}
\newmdenv[backgroundcolor=warncolor,linecolor=warnborder,linewidth=1.2pt,
  leftline=true,rightline=false,topline=false,bottomline=false,
  innerleftmargin=12pt,innerrightmargin=8pt,innertopmargin=8pt,
  innerbottommargin=8pt,skipabove=10pt,skipbelow=6pt]{warnbox}
\newmdenv[backgroundcolor=opencolor,linecolor=openborder,linewidth=1.2pt,
  leftline=true,rightline=false,topline=false,bottomline=false,
  innerleftmargin=12pt,innerrightmargin=8pt,innertopmargin=8pt,
  innerbottommargin=8pt,skipabove=10pt,skipbelow=6pt]{openproblem}

\theoremstyle{plain}
\newtheorem{theorem}{Theorem}[section]
\newtheorem{lemma}[theorem]{Lemma}
\newtheorem{proposition}[theorem]{Proposition}
\newtheorem{conjecture}[theorem]{Conjecture}
\theoremstyle{definition}
\newtheorem{definition}[theorem]{Definition}
\newtheorem{remark}[theorem]{Remark}

\newcommand{\ZSSP}{Z_{\mathrm{SSP}}^{*}}
\setlist[itemize]{noitemsep,topsep=4pt}
\setlist[enumerate]{noitemsep,topsep=4pt}

\begin{document}

\begin{center}
  {\color{navyblue}\rule{\textwidth}{2pt}}\\[0.4cm]
  {\LARGE\bfseries\color{navyblue}The Job Sequencing and Tool Switching Problem}\\[0.2cm]
  {\large\bfseries\color{navyblue}Part X: Position-Indexed Formulations}\\[0.1cm]
  {\large\itshape\color{midblue}Two new exact formulations with $O(n|T|)$ constraints}\\[0.3cm]
  {\color{navyblue}\rule{\textwidth}{2pt}}
\end{center}

\vspace{0.3cm}
\begin{mdframed}[backgroundcolor=lightblue,linecolor=navyblue,linewidth=1pt,
  innerleftmargin=14pt,innerrightmargin=14pt,innertopmargin=8pt,innerbottommargin=8pt]
\small\textbf{Scope and status.}
%% TODO-VERIFY: drafted by Claude-Fable 2026-06-10 from Shankar's variable
%% proposals (y_C^k; z_{C,C'}^k; hybrid). REWRITTEN same day after feedback:
%% each formulation is now STANDALONE (variables, constraints, objective,
%% exactness proof), independent of the solution machinery of Parts II/IV/VIII;
%% solution paradigms come after. All numerical claims:
%% _verification/verify_posform.py and verify_posform_f2.py.
This note introduces two \emph{standalone} exact MILP formulations of the uniform
SSP, both with constraint sets of size $O(n|T|)$ --- polynomial in the input, fixed in
advance --- while the (exponentially many) configuration variables may be handled by
column generation \emph{without ever adding or modifying a constraint}. They are, to
our knowledge, the first exact configuration-based SSP formulations with this
property; all known alternatives carry per-configuration ordering constraints (MTZ,
flow, SECs). \textbf{Proved}: exactness of both formulations; the plain LP relaxation
of the first is always $0$; tool-counting rows lift it to $\ge|U|-b$.
\textbf{Verified}: the second formulation's LP can \emph{strictly exceed} $|U|-b$
(the SSPMF root bound) and is tight on the $6$-ring. Solution paradigms (CG,
Benders) and open problems close the note.
\end{mdframed}

\vspace{0.3cm}
\tableofcontents
\newpage

% -----------------------------------------------------------------------------
\section{Motivation: Ordering Without Ordering Constraints}
\label{sec:goal}
% -----------------------------------------------------------------------------

Jobs $J=\{1,\dots,n\}$ with tool sets $T_j\subseteq T$, $|T_j|\le b\le|T|$;
$U=\bigcup_jT_j$; configurations $V=\{C\subseteq T:|C|=b\}$;
$H_j=\{C\in V:T_j\subseteq C\}$; transition cost $d(C,C')=|C'\setminus C|$.
Free-initial-load convention: $\ZSSP$ counts insertions after a free first
configuration.

Configuration-based formulations of the SSP must order the visited configurations.
Every classical ordering device prices this per \emph{configuration}: SECs are
exponentially many; MTZ needs a potential and rows for each configuration pair;
single-commodity flow needs a balance row per configuration. Since $|V|$ is
exponential, all of these tie the \emph{constraint} set to the configuration set ---
the obstruction identified in Part~IV.

\emph{Position indexing} (Shankar's proposal) removes the obstruction: time is
discretised into positions $k=1,\dots,n$, which are ordered \emph{a priori}, so no
ordering constraints exist at all. The price is that costs and feasibility must be
expressed through position-indexed variables; the two formulations below differ in
how they do this. Both rest on the following normal form.

\begin{thmbox}
\begin{lemma}[Trajectory normal form]
\label{lem:nf}
Every feasible instance has an optimal schedule whose magazine trajectory consists of
$m\le n$ \emph{pairwise distinct} configurations $C_{(1)},\dots,C_{(m)}$, each serving
at least one job, with cost $\sum_{l<m} d(C_{(l)},C_{(l+1)})$.
\end{lemma}
\end{thmbox}
\begin{proofbox}
\begin{proof}
Take any optimal schedule with full magazines (pad arbitrarily; standard). Its
trajectory has one configuration per job position. Collapse consecutive equal
configurations (zero cost change). If a configuration recurs non-consecutively,
reassign the jobs of the later visit to the earlier one (identical configuration)
and splice it out: by the triangle inequality $d(X,Y)\le d(X,C)+d(C,Y)$ (valid since
$Y\setminus X\subseteq(Y\setminus C)\cup(C\setminus X)$ for $|X|=|C|=|Y|=b$) the cost
does not increase. Delete job-less configurations by the same splice. Each step
shortens the sequence, so the process terminates; what remains has $m\le n$ since
each surviving configuration serves a job.
\end{proof}
\end{proofbox}

% -----------------------------------------------------------------------------
\section{Formulation PCF: Position--Configuration}
\label{sec:pcf}
% -----------------------------------------------------------------------------

\begin{formbox}
\textbf{PCF.} Variables: $y_C^k\in\{0,1\}$ ($C\in V$, $k=1,\dots,n$): configuration
$C$ occupies position $k$; $w_t^k\ge0$ ($t\in T$, $k=2,\dots,n$): insertion of tool
$t$ at position $k$. Abbreviate $a_t^k:=\sum_{C\ni t}y_C^k$.
\begin{align}
  \min\ & \sum_{t\in T}\sum_{k=2}^{n} w_t^k \notag\\
  \text{s.t.}\quad & \sum_{C\in V} y_C^k = 1 && k=1,\dots,n \label{eq:P}\\
  & \sum_{k=1}^{n}\sum_{C\in H_j} y_C^k \ \ge\ 1 && j\in J \label{eq:C}\\
  & w_t^k \ \ge\ a_t^k - a_t^{k-1} && t\in T,\ k=2,\dots,n \label{eq:W}\\
  & \sum_{k=2}^{n} w_t^k \ \ge\ 1 - a_t^1 && t\in U \label{eq:T}
\end{align}
Constraints: $2n+(n-1)|T|+|U| = O(n|T|)$. Variables: $n|V|+(n-1)|T|$.
\end{formbox}

The model says: one configuration per position (repeats allowed and free); every job
is covered at some position; an insertion of $t$ is charged whenever the presence
mass of $t$ increases from one position to the next; every used tool absent from the
first configuration is inserted at least once.

\begin{thmbox}
\begin{theorem}[PCF is exact]
\label{thm:pcf}
The optimal value of PCF with binary $y$ equals $\ZSSP$.
\end{theorem}
\end{thmbox}
\begin{proofbox}
\begin{proof}
($\le$) Take an optimal normal-form schedule (\cref{lem:nf}). Set $y^l_{C_{(l)}}=1$
for $l\le m$ and $y^k_{C_{(m)}}=1$ for $k>m$; set
$w_t^k=\max(0,a_t^k-a_t^{k-1})$. \eqref{eq:P}, \eqref{eq:C} hold by construction. At
binary $y$, $a_t^k\in\{0,1\}$ and $\sum_t\max(0,a_t^k-a_t^{k-1})=d$ between the
consecutive configurations, so the objective equals the schedule cost; repeats
contribute $0$. \eqref{eq:T}: a tool $t\in U$ with $a_t^1=0$ is required at some
covered position, hence its presence rises from $0$ to $1$ somewhere and
$\sum_kw_t^k\ge1$.

($\ge$) An integer-feasible solution reads off a trajectory $D_1,\dots,D_n$
(one configuration per position). Its insertions at step $k$ are exactly the tools
with $a_t^k-a_t^{k-1}=1$, each charged by \eqref{eq:W}, so the objective is at least
$\sum_k d(D_{k-1},D_k)$. By \eqref{eq:C} every job is covered by some $D_k$;
assigning each job to a covering position and ordering jobs by position yields a
feasible schedule of cost $\le$ the objective. Hence the value is $\ge\ZSSP$.
\end{proof}
\end{proofbox}

\begin{thmbox}
\begin{proposition}[Plain LP value is zero]
\label{prop:lp0}
Without \eqref{eq:T}, the LP relaxation of PCF is $0$ for every feasible instance.
\end{proposition}
\end{thmbox}
\begin{proofbox}
\begin{proof}
Pick $C_j\in H_j$ for each $j$ and set $y^k=\frac1n\sum_j\delta_{C_j}$ for
\emph{every} $k$, $w=0$. \eqref{eq:P} holds; \eqref{eq:C} gives
$n\cdot\frac1n=1$; $a_t^k$ is constant in $k$ so \eqref{eq:W} is slack at $w=0$.
(The position-symmetric analogue of the Tang--Denardo LP-zero pathology: blending one
mixture across all positions makes insertions, which are \emph{differences},
invisible.)
\end{proof}
\end{proofbox}

\begin{thmbox}
\begin{proposition}[Tool-counting rows are valid and give $\mathrm{LP}\ge|U|-b$]
\label{prop:T}
Rows \eqref{eq:T} hold at every integer solution, and the LP relaxation of the full
PCF is $\ge|U|-b$.
\end{proposition}
\end{thmbox}
\begin{proofbox}
\begin{proof}
Validity is the last step of the proof of \cref{thm:pcf}. For the bound, sum
\eqref{eq:T} over $t\in U$ and use
$\sum_{t}a_t^1=\sum_C|C|\,y_C^1=b$ from \eqref{eq:P}:
$\sum_{t,k}w_t^k\ge|U|-\sum_{t\in U}a_t^1\ge|U|-b$.
\end{proof}
\end{proofbox}

\begin{warnbox}
\textbf{Verified} (\texttt{verify\_posform.py}): IP $=\ZSSP$ on the $6$-ring,
$5$-ring and random instances; plain LP $=0.00$ throughout; LP with \eqref{eq:T}
$=|U|-b$ \emph{exactly} on every tested instance --- the counting rows are the entire
strength of the PCF relaxation so far. PCF thus matches the SSPMF root bound with a
configuration-variable formulation.
\end{warnbox}

% -----------------------------------------------------------------------------
\section{Formulation PTF: Position--Transition}
\label{sec:ptf}
% -----------------------------------------------------------------------------

The PCF objective sees only per-tool \emph{differences}; exact pair costs are
recovered by indexing variables by \emph{transitions} (Shankar's $y_{i,j}^k$).
Let $\bot\notin V$ be an absorbing terminal symbol and
$A=\{(C,C'):C\ne C'\in V\}\cup\{(C,\bot):C\in V\}\cup\{(\bot,\bot)\}$
(\emph{no diagonal pairs} $(C,C)$ --- this is essential).

\begin{formbox}
\textbf{PTF.} Variables $z^k_{(C,C')}\in\{0,1\}$ for $(C,C')\in A$, $k=1,\dots,n-1$:
the $k$-th step moves from $C$ to $C'$ ($C'=\bot$: the trajectory has ended).
Costs $d(C,C')$; $d(C,\bot)=d(\bot,\bot)=0$. Write
$\mathrm{hd}_t^k=\sum_{(C,C')\in A:\,t\in C'}z^k$,
$\mathrm{tl}_t^k=\sum_{(C,C')\in A:\,t\in C}z^k$ (with $t\in\bot$ false), and
$\mathrm{ins}_t^k=\sum_{(C,C')\in A:\,t\in C'\setminus C}z^k$.
\begin{align}
  \min\ & \sum_{k=1}^{n-1}\sum_{(C,C')\in A} d(C,C')\,z^k_{(C,C')} \notag\\
  \text{s.t.}\quad
  & \sum_{(C,C')\in A} z^k_{(C,C')} = 1 && k=1,\dots,n-1 \label{eq:Pz}\\
  & \mathrm{hd}_t^k \ =\ \mathrm{tl}_t^{k+1} && t\in T,\ k\le n-2 \label{eq:cons}\\
  & \textstyle\sum_{(C,C')\in A:\,C'=\bot} z^k_{(C,C')} \ \le\ z^{k+1}_{(\bot,\bot)}
    && k\le n-2 \label{eq:abs}\\
  & \textstyle\sum_{(C,C')\in A:\,T_j\subseteq C} z^1
    \;+\;\sum_{k}\sum_{(C,C')\in A:\,T_j\subseteq C'} z^k \ \ge\ 1 && j\in J
    \label{eq:Cz}\\
  & \sum_{k=1}^{n-1}\mathrm{ins}_t^k \ \ge\ 1-\mathrm{tl}_t^1 && t\in U \label{eq:Tz}
\end{align}
Constraints: $(n-1)+(n-2)|T|+(n-2)+n+|U|=O(n|T|)$. Variables: $O(n|V|^2)$.
\end{formbox}

\eqref{eq:cons} makes consecutive steps chain (the tools of step $k$'s head are the
tools of step $k{+}1$'s tail --- aggregated \emph{per tool}, which keeps the rows
polynomial); \eqref{eq:abs} makes $\bot$ absorbing; \eqref{eq:Cz} covers each job at
the first tail or any head; \eqref{eq:Tz} are the counting rows.
%% TODO-VERIFY(Claude-Fable 2026-06-10, self-review): n>=2 assumption made
%% explicit — with n=1 there are no steps and the coverage row is vacuously
%% infeasible; the n=1 instance is trivial (Z*=0, any covering configuration).
Throughout assume $n\ge2$ (for $n=1$ the SSP is trivial and PTF has no steps).
At integer points \eqref{eq:cons} alone already forces absorption ($\bot$ has no
tools, and a configuration \emph{is} its tool set); \eqref{eq:abs} is kept because
it tightens the LP. Because diagonal
pairs are excluded, a step either advances to a \emph{different} configuration
(cost $\ge1$) or terminates --- blending cannot hide in free self-loops.

\begin{thmbox}
\begin{theorem}[PTF is exact]
\label{thm:ptf}
The optimal value of PTF with binary $z$ equals $\ZSSP$.
%% TODO-VERIFY(Claude-Fable): the (>=) direction below relies on per-tool
%% consistency forcing consecutive steps to chain through the SAME configuration
%% at integer points: hd and tl are 0/1 tool-indicator vectors and a config is
%% determined by its tool set, so equality of all |T| coordinates identifies the
%% configuration. Check the boundary case where step k is absorbed (head = bot).
\end{theorem}
\end{thmbox}
\begin{proofbox}
\begin{proof}
($\le$) Embed a normal-form trajectory (\cref{lem:nf}): set
$z^l_{(C_{(l)},C_{(l+1)})}=1$ for $l<m$, $z^{m\dots n-1}$ on $(C_{(m)},\bot)$ then
$(\bot,\bot)$. All rows hold (consecutive distinctness of the normal form is what
the diagonal-free $A$ requires); the cost telescopes to the schedule cost;
\eqref{eq:Tz} holds as in \cref{thm:pcf}.

($\ge$) At an integer solution, exactly one pair is chosen per step. By
\eqref{eq:cons} the head tool-vector of step $k$ equals the tail tool-vector of step
$k{+}1$; since a configuration is its tool set, the chosen pairs chain into a
trajectory $D_1\to D_2\to\cdots$ until absorption, after which \eqref{eq:abs} keeps
all later steps at $\bot$. The objective is exactly $\sum d(D_l,D_{l+1})$. By
\eqref{eq:Cz} every job is covered by a visited configuration; assign and order as in
\cref{thm:pcf}. Hence the value is $\ge\ZSSP$.
\end{proof}
\end{proofbox}

\begin{warnbox}
\textbf{Verified} (\texttt{verify\_posform\_f2.py}): IP $=\ZSSP$; and the LP
relaxation is strictly stronger than PCF's:
\begin{center}
\renewcommand{\arraystretch}{1.2}
\begin{tabular}{lccccc}
\toprule
instance & $\ZSSP$ & LP(PTF) & LP(PTF, no \eqref{eq:Tz}) & LP(PCF) & $|U|-b$\\
\midrule
$6$-ring ($b{=}3$) & 3 & \textbf{3.00} & 2.00 & 3.00 & 3\\
above-bounds ($b{=}3$, $\ZSSP>\max$) & 3 & \textbf{2.10} & 2.05 & 2.00 & 2\\
\bottomrule
\end{tabular}
\end{center}
On the $6$-ring the PTF LP is \emph{tight}; on an instance whose optimum strictly
exceeds $\max(K^*{-}1,|U|-b)$, the PTF LP \textbf{strictly exceeds $|U|-b$} ---
something neither PCF nor (at the root) SSPMF's bound provides. Even without the
counting rows the PTF LP is strictly positive (no Tang--Denardo collapse): with
diagonals allowed it would be $0$ (verified) --- excluding them is what buys the
strength.
\end{warnbox}

\begin{remark}[Relation to the hybrid proposal $y_C^k$ $+$ $x_{C,C'}$]
\label{rem:hybrid}
The natural hybrid keeps PCF's $y_C^k$ and adds configuration-arc variables
$x_{C,C'}$ carrying the cost, linked by $x_{C,C'}\ge y_C^k+y_{C'}^{k+1}-1$. Exact
linking requires one row per (pair, position) --- the row set again grows with the
generated configurations, violating the design criterion; aggregating the linking
per tool reproduces exactly \eqref{eq:cons}. In this precise sense \textbf{PTF
\emph{is} the hybrid formulation}: $z^k_{(C,C')}$ realises $x_{C,C'}$
position-wise, and the tails/heads play the role of $y_C^k$.
\end{remark}

% -----------------------------------------------------------------------------
\section{Can Per-Configuration MTZ Be Row-Reduced?}
\label{sec:mtzreduce}
% -----------------------------------------------------------------------------

Shankar asked whether the per-configuration MTZ of Part~IV can instead be thinned to
polynomially many rows. The obstacle is structural: MTZ rows
$v_C-v_{C'}+Nx_{C,C'}\le N-1$ are indexed by \emph{arc variables}, and any ordering
certificate (potentials, single-commodity flow balances, GG-flows) needs at least
one row per configuration that can appear --- the certificate must ``know'' each
node. Aggregating potentials over clusters of configurations is exactly the
cluster-MTZ of Part~IV, which is provably not exact. Position indexing sidesteps the
issue by making order \emph{structural} rather than certified; we therefore regard
PCF/PTF as the row-reduction of per-configuration MTZ, and pose the converse as
\cref{op:x4}: a proof that no exact configuration-arc formulation has
$o(|V_{\mathrm{used}}|)$ ordering rows.

% -----------------------------------------------------------------------------
\section{Comparison and Solution Paradigms}
\label{sec:verdict}
% -----------------------------------------------------------------------------

\begin{center}
\renewcommand{\arraystretch}{1.3}
\begin{tabular}{lcccc}
\toprule
formulation & rows & IP exact & root LP & columns\\
\midrule
per-config MTZ (Part IV) & grows with $V_{\mathrm{gen}}$ & yes & untested & $O(|V|^2)$\\
cluster MTZ (Part IV) & $O(n^2)$ & \textbf{no} & weak & $O(|V|^2)$\\
\textbf{PCF} & $O(n|T|)$, fixed & \textbf{yes} & $=|U|-b$ (verified) & $n|V|$\\
\textbf{PTF} & $O(n|T|)$, fixed & \textbf{yes} & $\ge$ PCF; $>|U|-b$ possible & $O(n|V|^2)$\\
\bottomrule
\end{tabular}
\end{center}

\textbf{Solution paradigms} (consequences, not design drivers). \emph{Column
generation}: in both formulations a generated configuration (PCF) or configuration
pair (PTF) contributes columns to existing rows only; pricing is a set-selection
problem over $b$ tools (PCF: single set with job-coverage bonuses, the JGP pricing
class; PTF: a coupled pair --- harder, and the practical price of its stronger LP).
\emph{Benders}: in PTF one may project the $z$-layer onto a surrogate $\theta$ per
position and cut from the per-tool consistency duals; whether this beats direct CG
is untested. \emph{Hybrids with the job-level BBC} (Part~VIII) are an implementation
topic deliberately kept out of this note.

\begin{openproblem}
\textbf{OP-X1.} Characterise the PTF LP bound; in particular prove (or refute)
$\mathrm{LP(PCF)}\le\mathrm{LP(PTF)}$ in general, quantify how far past $|U|-b$ the
PTF LP can go (it reached $\ZSSP$ on the ring), and compare against the Catanzaro
F3/F4 root bounds on Tabela1C.
\end{openproblem}
\begin{openproblem}
\textbf{OP-X2.} Polynomial row families beyond the counting rows: position-conflict
cliques, per-tool interval/spread rows tying $w_t$ to the positions of jobs needing
$t$, and symmetry breaking (reversal + position shifts); which are facet-defining for
the PCF/PTF polytopes on small instances (PORTA)?
\end{openproblem}
\begin{openproblem}
\textbf{OP-X3.} Practical CG for PTF: pair pricing as a quadratic set-selection
(two coupled $b$-subsets with $d$-cost and dual rewards) --- exact MILP pricing
vs.\ heuristic pricing with exact fallback; degeneracy management for the $n-1$
identical step rows.
\end{openproblem}
\begin{openproblem}
\label{op:x4}
\textbf{OP-X4.} Lower bound for arc-based ordering: prove that every exact
formulation of the SSP whose integer solutions are characterised by
configuration-arc incidence requires $\Omega(|V_{\mathrm{used}}|)$ constraints, or
exhibit a sub-linear ordering certificate (this would also row-reduce per-config
MTZ directly, Shankar's original question).
\end{openproblem}

\bibliographystyle{abbrvnat}
\bibliography{references}
\end{document}
